\documentclass[UTF8,a4paper,12pt]{ctexart}

\usepackage{geometry}
\geometry{left=2.6cm,right=2.6cm,top=2.6cm,bottom=2.6cm}
\usepackage{amsmath,amssymb}
\usepackage{booktabs,tabularx,array,makecell,multirow}
\usepackage{graphicx}
\usepackage{pdflscape}
\usepackage{float}
\usepackage{caption}
\usepackage{enumitem}
\usepackage[round,authoryear]{natbib}
\usepackage[hidelinks]{hyperref}
\usepackage{xcolor}
\usepackage[normalem]{ulem}
\usepackage{graphicx}
\usepackage{svg}
\captionsetup[table]{name=表}
\numberwithin{table}{section}
\renewcommand{\thetable}{\thesection-\arabic{table}}
\renewcommand{\arraystretch}{1.18}
\setlength{\parindent}{2em}
\setlength{\parskip}{0.25em}
\sloppy

\title{\heiti 面向未见概念归纳知识追踪的目标条件认知证据检索\\[0.5em]
\large Target-Conditioned Cognitive Evidence Retrieval for Inductive Knowledge Tracing over Unseen Concepts}
\author{}
\date{}

\begin{document}
\maketitle

\begin{abstract}


在线教育平台中的题目和知识概念持续更新，测试目标可能关联训练阶段未见概念，使依赖题目或概念标识的知识追踪模型缺少直接监督信号。本文提出 Cognitive-RAG KT：以目标题目为查询，从学生目标前历史中检索可迁移认知证据，结合题目语义、知识结构和协同行为构造分层认知上下文；同时用结构化认知摘要压缩掌握点、薄弱点、迁移状态和证据质量。随后，模型将学生侧上下文、摘要派生特征和目标题目先验映射为 soft slots，注入冻结 LLM，并由标签 logits 输出作答正确概率。在 MOOCRadar、XES3G5M 和 FoundationalASSIST 上，Cognitive-RAG KT 优于多类 KT 基线；切片和消融结果表明，目标条件认知证据检索提供最稳定的概率排序增益，结构化证据和结构化认知摘要分别提供补充支持。
\end{abstract}

\noindent\textbf{关键词：}知识追踪；未见概念归纳预测；目标条件认知证据检索；结构化认知摘要；冻结 LLM

\section{引言}



知识追踪（Knowledge Tracing, KT）旨在依据学生既往学习交互估计其知识状态，并预测其在后续题目或学习活动中的表现。作为智能教育系统的基础能力，KT 为个性化练习推荐、学习路径规划和教学干预提供依据。然而，真实教育平台中的题库、课程内容与知识组织并非静态不变：新的题目和知识概念会持续出现，而这些概念可能不包含在模型训练阶段的监督交互中。对于关联未见知识概念的目标题目，模型无法直接依赖该概念的历史标签、题目 ID 统计或概念嵌入完成预测，因而面临从既有学习经验迁移至新概念目标的归纳难题。



多数传统 KT 方法以题目 ID、知识概念 ID 或其监督学习表示为核心输入。BKT、DKT 和 DKVMN 分别通过概率图、循环网络或记忆网络建模学生历史交互，SAKT、SAINT 和 AKT 则进一步利用注意力机制捕捉交互依赖\citep{corbett1995knowledge,piech2015deep,zhang2017dynamic,pandey2019sakt,choi2020saint,ghosh2020akt}。这些方法在训练与测试共享相同题目或概念分布时能够有效学习学生状态；但当测试目标关联未见知识概念时，离散标识及其参数化表示缺少直接学习依据。已有研究开始利用题目文本、概念语义和结构关系缓解该问题\citep{fu2024sinkt}，但仍需回答一个更细粒度的问题：面对同一未见概念目标，不同学生历史中的哪些交互真正具有迁移价值。



未见概念预测不能通过简单扩展历史长度解决。学生历史中同时包含与当前目标密切相关的前置支撑、语义近邻和结构关联交互，也包含弱相关或重复信息；同一条历史交互面对不同未见概念目标具有不同的迁移价值。因此，关键问题不是向模型输入更多历史，而是判断哪些目标前交互能够为当前未见概念目标提供可迁移证据。

为此，本文提出 Cognitive-RAG KT。模型以当前未见概念目标题目为查询，从学生目标前历史中执行目标条件认知证据检索，并结合题目语义、知识概念、认知属性、知识结构和协同行为关系构造分层认知上下文。结构化认知摘要进一步把检索证据压缩为掌握点、薄弱点、迁移状态、风险水平和证据质量等受约束字段；该摘要不生成预测标签，而是作为学生侧上下文特征参与概率判别。随后，学生侧上下文、摘要派生特征和目标题目侧先验被映射为连续 soft slots，注入冻结 LLM 的表示空间，并通过标签 logits 的二分类 softmax 输出作答正确概率。

因此，本文将检索增强适配为跨概念迁移中的目标条件证据选择，而不是把完整历史转写为自然语言 prompt 或让 LLM 自由生成判断。实验结果显示，本文方法在三个公开数据集上取得稳定优势；消融分析表明，目标条件认知证据检索贡献最大，结构化证据提供次级支持，结构化认知摘要提供辅助压缩增益。


本文的贡献包括四个方面。第一，本文构建未见概念归纳知识追踪评估协议，并在该协议下研究目标条件历史证据选择对学生作答预测的作用；测试目标至少关联训练阶段未见的知识概念，模型需要从学生在既有概念上的学习经历迁移知识。第二，本文提出目标条件认知证据检索机制，以当前未见概念目标题目为查询，从学生目标前历史中选择具有迁移价值的交互，并结合题目语义、知识结构和协同行为关系构造分层认知上下文，同时通过结构化认知摘要对检索证据中的掌握点、薄弱点、迁移状态和证据质量进行紧凑压缩。第三，本文设计基于冻结 LLM 的 soft-slot 概率预测模块，将学生侧上下文、摘要派生特征与目标题目侧先验映射为连续条件输入，并通过标签 logits 而非自由文本生成完成作答正确概率估计。第四，本文在 MOOCRadar~\citep{yu2023moocradar}、XES3G5M~\citep{liu2023xes3g5m} 和 FoundationalASSIST~\citep{worden2026foundationalassist} 三个公开数据集的未见概念归纳评估协议下，对比多类代表性知识追踪基线模型，并通过检索对照和关键组件消融分析所提框架的有效性。

\section{相关工作}

\subsection{传统知识追踪与未见知识概念预测}



知识追踪（Knowledge Tracing, KT）通常根据学生历史学习行为估计其知识状态，并预测其在后续题目上的作答表现。早期模型以 BKT、项目反应理论和认知诊断模型为代表，通过显式技能、题目参数或隐变量描述学生状态\citep{corbett1995knowledge,chen2021irt,pavlik2009pfa,tatsuoka1983rulespace}。随后，DKT 和 DKVMN 使用循环或记忆结构学习交互依赖，SAKT、SAINT 和 AKT 则利用注意力机制建模更复杂的长程关系\citep{piech2015deep,zhang2017dynamic,pandey2019sakt,choi2020saint,ghosh2020akt}。这些方法共同推动了 KT 的预测能力，但其有效性通常建立在训练和测试共享相同或高度重叠的题目、概念及其监督分布之上。



当测试目标题目关联的知识概念未进入训练监督时，上述共享分布假设不再成立。模型无法从该目标概念的历史标签中学习稳定表示，也难以仅凭离散题目或概念标识推断其与既有学习经历之间的关系。已有工作尝试利用题目文本、概念语义和题目--概念结构提升模型对新概念的表征能力\citep{fu2024sinkt}。这说明未见概念预测的关键不是重复估计已见概念上的常规时序模式，而是建立从既有概念经验到目标新概念的可迁移联系。









然而，静态概念关系不足以确定学生在当前目标上的有效证据。同一未见概念对不同学生可能需要不同历史支持；同一学生面对不同未见目标也会激活不同学习片段。深度 KT 的注意力机制虽可突出部分历史，但多为隐式加权，难以说明哪些交互承担跨概念迁移作用。本文因此关注目标条件下的个体化证据选择：在训练阶段概念监督缺失的约束下，从学生历史中检索支持当前目标预测的可追溯认知证据。

本文沿用 inductive KT 中关于训练与测试题目/概念分布不完全重合的设定，将关注点进一步限定到测试目标包含训练阶段未见知识概念的场景；为避免与新学生冷启动或新题目冷启动混淆，后文将其称为未见概念归纳预测。

\subsection{检索增强与个体化认知证据建模}



检索增强生成（Retrieval-Augmented Generation, RAG）的核心思想不是简单扩展模型输入，而是根据当前查询从候选信息中选择相关证据，并以这些证据补充或约束后续推断。Lewis 等提出的 RAG 将参数化语言模型与非参数化外部记忆结合，在生成答案前检索相关文本，从而改善知识密集型任务中的信息访问能力\citep{lewis2020rag}。后续工作进一步表明，检索器、重排器、上下文选择与证据组织方式会共同影响语言模型输出的可靠性。



在教育场景中，检索增强已被用于课程问答、学习资源推荐、自动反馈生成和知识库访问等任务\citep{li2025ragedusurvey,swacha2025ragchatbots,dong2023kgrag,zheng2024edukdqa}。这些研究说明，检索能够将模型的输出锚定在可检查的外部证据上。然而，知识追踪并不需要从外部文档中生成文本答案，而需要依据学生先前行为预测其对当前题目的作答概率。因此，适用于 KT 的检索对象应从外部文档转变为学生自身的目标前学习历史，检索目标也应从回答开放问题转变为识别支持当前预测的个体化认知证据。



对于未见概念预测，检索进一步承担跨概念迁移的职责。当前目标题目作为查询，学生目标前历史作为候选证据池；题目文本、概念语义、认知属性、知识结构和协同行为关系用于评估候选历史与当前目标之间的迁移关系。检索输出不是未经筛选的交互列表，而是具有支持方向、认知角色、相对位置和关系来源的结构化证据集合。这样的证据既减少弱相关历史占据 LLM 上下文的风险，也使后续概率判断能够追溯到具体的学生学习经历。

\subsection{大语言模型的推理潜力及其在知识追踪中的约束使用}



大语言模型在文本理解、语义关联和条件推断方面展现出较强潜力，已被用于教育问答、学习资源生成、个性化反馈和学习分析等任务\citep{kasneci2023chatgpt,cho2024systematic}。对于未见知识概念，题目文本和概念描述为 LLM 提供了理解目标能力要求的语义入口；当模型同时接收学生既有学习状态与目标概念的结构化描述时，LLM 可以发挥整合异构证据和进行条件判别的作用。



已有 LLM-enhanced KT 方法尝试使用语言模型编码题目与概念语义，或通过指令对齐和上下文表示辅助下一题预测\citep{fu2024sinkt,wang2025llmkt}。这些研究为 LLM 参与概念泛化提供了基础，但 LLM 的一般语言知识不等同于学生个体的知识状态。若直接依据题目文本或自由生成输出进行预测，模型可能受到提示格式、外部常识、生成随机性和后处理规则影响，难以保证判断确实建立在学生历史表现之上。



本文因此不将 LLM 用作开放式推理主体，而是将其限定为冻结的条件表示与标签判别空间。目标条件认知证据检索首先为 LLM 提供与当前目标相关、可追溯且结构化的学生证据；结构化认知摘要仅将这些证据压缩为受约束摘要字段，不承担最终标签预测；学生侧证据状态、摘要派生特征和目标题目侧先验随后映射为连续 soft slots，并与标签约束提示共同输入模型。预测阶段只读取正确与错误标签的 logits 并计算二分类概率。该设计保留 LLM 对异构语义证据进行条件整合的潜力，同时限制其依据无关常识或自由生成结果作出判断。





综上，本文不以检索增强替代既有 KT，也不将 LLM 视为脱离学生历史的知识来源。本文利用目标条件认知证据检索，将学生在既有概念上的学习经历转化为面向未见概念目标的可迁移证据；再由分层认知上下文和结构化认知摘要形成学生侧条件表示，并通过 soft slots 约束性注入冻结 LLM，最终从标签 logits 计算作答正确概率。由此，检索限定证据边界，上下文与摘要提供可编码的学生侧表示，冻结 LLM 提供受约束的条件判别空间，三者共同服务于未见知识概念上的归纳预测。

\section{方法}

\subsection{研究问题、信息边界与总体框架}

\noindent\textbf{未见知识概念上的归纳预测}



设学生集合、题目集合和知识概念集合分别为 \(\mathcal{U}\)、\(\mathcal{Q}\) 和 \(\mathcal{C}\)。对于任一学生 \(u\in\mathcal{U}\)，其学习过程由按时间顺序排列的交互序列组成：
\[
S_u=\{(q_{u,1},y_{u,1}),(q_{u,2},y_{u,2}),\ldots,(q_{u,T_u},y_{u,T_u})\}
\]



其中，\(q_{u,i}\in\mathcal{Q}\) 表示学生 \(u\) 在第 \(i\) 次交互中作答的题目，\(y_{u,i}\in\{0,1\}\) 表示对应作答结果。对于第 \(t\) 个预测位置，当前需要预测的题目记为目标题目 \(q_{u,t}\)，模型只能使用学生在该目标之前发生的历史交互：
\[
H_{u,<t}=\{(q_{u,i},y_{u,i})\mid 1\leq i<t\}
\]

本文将知识追踪预测定义为如下条件概率估计问题：
\[
p_{u,t}=P(y_{u,t}=1\mid H_{u,<t},q_{u,t})
\]



其中，\(p_{u,t}\) 表示学生 \(u\) 在目标题目 \(q_{u,t}\) 上作答正确的理论目标概率，\(\hat{p}_{u,t}\) 表示模型对该概率的估计。训练阶段仅使用真实标签 \(y_{u,t}\) 计算监督损失。预测阶段，目标标签不进入模型输入。

对于未见概念目标，预测还依赖于题目在作答前可获得的属性 \(a(q_{u,t})\)。这些属性包括题干文本、人工标注概念和静态认知属性。后文的多源先验模块据此将一般形式的条件概率细化为 \(P(y_{u,t}=1\mid H_{u,<t},q_{u,t},a(q_{u,t}))\)。

为定义未见概念归纳预测的评估边界，设用于模型拟合和测试的知识概念集合分别为 
\(\mathcal{C}_{\mathrm{train}}\) 和 \(\mathcal{C}_{\mathrm{test}}\)，
且 \(\mathcal{C}_{\mathrm{train}}\cap\mathcal{C}_{\mathrm{test}}=\emptyset\)。对于测试目标 \(q_{u,t}\)，其关联概念集合满足 \(C(q_{u,t})\cap\mathcal{C}_{\mathrm{test}}\neq\emptyset\)。因此，模型拟合阶段未接收与目标概念相关的监督交互。模型必须依据学生在既有概念上的学习经历和目标题目可获得先验完成迁移预测。

本文将该设定称为训练阶段未见概念上的归纳预测。该设定不额外假设目标前历史中完全不存在同测试概念交互。若采用更严格的历史遮蔽设置，应在实验中单独报告。



对于未见概念目标，学生历史中的交互不只具有一般时序意义，还可能承担跨概念迁移作用。同一条历史记录面对不同目标时可能是前置支撑、语义近邻、同质迁移或结构补充，也可能与当前目标弱相关。因此，本文不将 \(H_{u,<t}\) 压缩为与目标无关的固定状态，而是将其视为围绕当前未见目标题目进行证据选择的候选历史集合。

\noindent\textbf{题目先验与可迁移证据表示}



在未见概念条件下，题目不能仅由离散题目编号表示，还必须由作答前可获得的认知属性描述。对于任一题目 \(q\)，本文从题目文本、关联概念、认知层级、难度属性、语义表示、图结构关系和协同行为关系等角度刻画其认知特征。

题目文本和人工概念标注为未见目标提供语义入口。认知层级和题型结构描述其能力要求。图结构与协同行为关系则建立该目标与既有概念、既有题目之间的可迁移连接。上述属性共同构成题目侧多源认知先验。

对于目标题目 \(q_{u,t}\)，题目侧认知先验用于描述当前预测目标的知识要求和认知要求。对于历史题目 \(q_{u,i}\)，题目侧认知先验用于判断该历史题目与目标题目之间的关系。

本文特别区分题目的原始编号与题目的语义标识表示。分层语义标识用于刻画题目在语义空间中的层次位置，不等同于原始题号。对于历史交互 \((q_{u,i},y_{u,i})\)，本文将其视为一个候选认知证据。该证据包含三类信息。第一类是历史题目自身的认知属性。第二类是学生在该历史题目上的作答结果。第三类是历史题目与当前目标题目之间的关系。

历史作答结果说明学生在该题上的表现，但其预测意义取决于该历史题目与目标题目的关系。例如，学生答对一个与目标题目共享知识概念的历史题，通常比答对一个无关题目更能支持当前预测。学生答错一个前置题，也可能提示其在当前目标题目上存在风险。

因此，本文将历史交互的证据价值定义为目标条件下的关系属性，而不是题目本身的静态标签。每条候选证据在当前预测中具有特定的认知角色，例如前置支撑、同质迁移、语义近邻、高阶迁移、图结构补充或协同补充。认知角色用于说明该历史交互为何与当前目标题目相关，并为后续证据检索和上下文构造提供结构化依据。

\noindent\textbf{目标前历史与预测时信息边界}



本文同时遵循目标前历史约束和未见概念归纳预测的信息边界。对于目标样本 \((u,t,q_{u,t})\)，模型只能使用 \(t\) 之前的学生历史交互 \(H_{u,<t}\) 与目标题目作答前可获得的属性。模型不使用目标标签 \(y_{u,t}\)，也不使用未来交互。目标标签仅在训练阶段用于计算损失，在评价阶段用于计算指标。该边界贯穿目标条件认知证据检索、分层认知上下文构造、结构化认知摘要生成和 soft-slot 概率预测四个环节。



除题干文本、人工标注概念和认知层级等静态属性外，题目难度、协同行为关系、知识图统计和行为嵌入均依赖学生作答记录。为避免将测试概念信息经由行为统计泄漏到模型中，这类先验仅由拟合集交互构建，并在验证与测试阶段保持冻结。任一验证或测试目标的标签、未来交互以及验证/测试集标签均不得参与题目先验构建、证据选择或模型输入。

首先，目标条件认知证据检索模块只从 \(H_{u,<t}\) 中选择历史证据，检索对象不包含目标题目作答结果和未来交互。其次，分层认知上下文构造模块只汇总目标前历史证据及其结构化关系，将其组织为学生侧认知上下文。第三，结构化认知摘要生成模块仅压缩已筛选的目标前证据及其统计摘要，不接收目标标签，也不输出最终预测标签。

最后，基于冻结 LLM 的 soft-slot 概率预测模块只接收学生侧认知上下文、结构化认知摘要派生特征、目标题目侧认知先验和标签约束提示，并输出作答正确概率。

通过上述定义，本文将知识追踪转化为目标条件认知证据选择与概率预测问题。模型输入不是完整历史拼接，而是围绕当前目标题目筛选和组织后的目标前认知证据；模型输出不是自然语言解释或开放式答案，而是作答正确预测概率 \(\hat{p}_{u,t}\)。这一定义为题目侧多源认知先验建模、目标条件认知证据检索、分层认知上下文构造、结构化认知摘要生成和 soft-slot 概率预测提供统一基础。

\noindent\textbf{面向未见知识概念的总体框架}



给定目标样本 \((u,t,q_{u,t})\)，Cognitive-RAG KT 以目标前历史 \(H_{u,<t}\) 和目标题目在作答前可获得的属性为输入。模型输出作答正确概率 \(\hat{p}_{u,t}\)。

当 \(C(q_{u,t})\cap\mathcal{C}_{\mathrm{test}}\neq\emptyset\) 时，模型在拟合阶段未获得该概念的监督交互。因此，框架的关键不在于为新概念学习新的离散标识，而在于把目标题目侧先验转化为连接该目标与既有学习经历的迁移接口。

框架依次执行五个步骤：题目侧多源认知先验建模、目标条件认知证据检索、分层认知上下文构造、结构化认知摘要生成，以及基于冻结 LLM 的 soft-slot 概率预测。先验和检索模块依据目标题目属性建立跨概念关系，并从既有概念历史中筛选可迁移证据。上下文和摘要模块将目标条件证据组织为学生侧条件表示。预测模块再把学生侧表示与目标先验编码为连续 soft slots，从而避免将完整历史压缩为与目标无关的固定状态。

\begin{figure}[t]
    \centering
    \includesvg[inkscapelatex=false,width=\textwidth]{figures/model}
    \caption{Cognitive-RAG KT 的总体框架。模型以目标题目和学生目标前历史为输入，基于目标题目条件从历史窗口中检索 Top-K 可迁移认知证据；随后将证据组织为分层认知上下文，并生成结构化认知摘要以压缩掌握点、薄弱点、迁移状态、风险水平和证据质量。学生侧上下文与目标题目侧先验分别映射为 context/target soft slots，注入冻结 LLM，并通过标签 logits 输出作答正确概率。}
    \label{fig:cognitive-rag-overview}
\end{figure}



在信息流上，目标题目侧先验包含文本语义、分层语义标识、认知增强表示、协同关系和知识结构信息。其中，文本及人工标注属性为未见概念提供直接描述。关系属性用于识别其与既有概念之间可能成立的迁移路径。

目标条件认知证据检索模块据此从学生目标前历史中召回并精排与该未见目标相关的交互。每条保留证据包括历史题目、作答结果、历史位置、证据角色和支持分数。分层认知上下文模块将这些离散证据组织为全局认知状态、局部证据块和证据结构特征。结构化认知摘要进一步补充掌握点、薄弱点、迁移状态和证据质量等压缩信息。

学生侧表示与目标题目侧先验随后分别映射为 context 与 target soft slots。两类 slots 被注入受标签约束的提示模板。冻结 LLM 仅读取候选标签 logits，并通过 softmax 计算 \(\hat{p}_{u,t}\)。

\subsection{多源认知先验与目标条件认知证据检索}

\noindent\textbf{可迁移的多源认知先验}



未见概念的核心困难在于：离散题目或概念 ID 只在拟合数据中具有可学习的参数意义，而测试目标的 ID 不提供可直接复用的监督信号。为此，本文将题目视为由文本内容、关联概念、认知属性及其与既有学习对象的关系共同刻画的对象，并构建题目侧多源认知先验。

对于题目集合 \(\mathcal{Q}=\{q_1,q_2,\ldots,q_M\}\)，每个题目 \(q_j\) 由题干文本、关联概念、认知属性以及可在未见概念评估边界内估计的统计特征描述。该表示不是对新概念 ID 的替代记忆，而是使模型能够基于可观察属性判断新目标与历史题目之间是否存在可迁移关系的共同坐标。

多源认知先验建模模块的目标，是为每个题目生成面向检索和预测的先验表示 \(r_j\)，使其既能刻画题目自身含义，又能反映题目在学习行为网络中的位置。首先，本文构建题目文本语义先验。题目文本包含了问题表述、解题对象、操作要求以及隐含的知识考查方式，是判断两个题目是否具有内容相关性的基础。本文将题目文本映射到语义向量空间，得到文本语义表示 \(s_j\)。

在该表示空间中，语义相近的题目具有更接近的向量表示，后续检索模块据此识别与目标题目在题干含义、知识表达或解题要求上相近的历史题目。与仅依赖题目 ID 或知识点标签相比，文本语义先验能够补充知识点标注粒度不足带来的信息缺失，使检索过程不仅关注是否属于同一知识点，也关注是否以相近方式考查相关能力。

其次，本文引入分层语义标识先验。分层语义标识并不是原始题号，而是对题目在语义层级或题目簇结构中的位置进行编码。该先验用于刻画题目之间由粗到细的层级关系。较高层级反映题目所属的宏观语义类别或知识区域，较低层级反映更细粒度的题目相似性。

对于题目 \(q_j\)，本文使用分层语义标识表示 \(h_j\) 描述其在层级语义空间中的位置。该表示使后续检索能够同时考虑细粒度近邻题目和较高层次上的相关题目，避免只在完全相同知识点或局部文本相似题目中选择证据。

在文本语义与分层语义标识之外，本文进一步构建认知增强语义先验。题目之间的关系不仅取决于文本相似度，还受到题目难度、认知要求、题型结构和知识组织方式的影响。两个文本相近的题目如果难度差异较大，其对学生当前掌握状态的证据意义并不相同。两个文本差异较大的题目如果共享相似的认知操作或解题结构，也可能提供有效的迁移证据。

因此，本文将题目难度、认知方向或结构性信息融入语义表示，得到认知增强语义表示 \(c_j\)。该表示用于刻画题目在认知层面的接近性，使检索模块能够识别同一能力要求下的相关历史交互，而不是仅依赖表层文本相似。

\noindent\textbf{难度、认知层级与状态特征}

题目侧先验描述的是题目本身及其与其他题目的关系，但知识追踪预测还需要结合学生在历史交互中的作答状态。对于同一道历史题目，学生答对与答错所传递的认知含义不同。对于同样的作答结果，其证据价值还会受到该题难度、历史位置以及学生近期表现的影响。

为此，本文在题目先验之外构建动态学生状态特征，用于解释历史作答结果在当前预测时刻下的含义。设学生 \(u\) 在目标题目 \(q_{u,t}\) 之前的历史交互为 \(H_{u,<t}\)。

对于任一历史交互 \((q_i,y_i)\)，其中 \(i<t\)，题目 \(q_i\) 的难度和认知层级为该交互提供题目侧背景，作答标签 \(y_i\) 反映学生在该题上的实际表现。当学生在低难度题目上连续答错时，该历史片段更可能反映基础概念掌握不足。当学生在高难度或综合性题目上答错时，该结果可能指向高阶迁移能力或复杂问题解决能力不足。相反，学生在与目标题目相关的高难度题目上答对，通常比在简单题目上答对提供更强的正向证据。

本文通过将题目难度、认知层级和历史作答结果共同纳入状态特征，使候选证据不再只是某题是否答对的离散记录，而是具有认知解释的历史状态片段。动态学生状态特征还用于描述学生在目标前历史中的近期趋势和稳定性。

知识掌握状态具有时间演化特征，同一学生早期的错误和近期的正确作答对当前预测的影响不同。本文在构造学生侧状态时保留历史位置、近期表现和局部作答模式，使后续检索与上下文构造能够区分长期背景证据和近期状态证据。

这样，历史交互在进入证据集合之前已经具有两类信息：一类来自题目本身的多源认知先验，另一类来自学生个体在该题上的实际作答状态。二者共同决定该历史交互对当前目标题目的支持方向和支持强度。

基于上述建模，本文将题目侧多源认知先验与学生侧动态状态特征结合，形成面向检索的候选证据表示。对于历史交互 \((q_i,y_i)\)，其证据表示由题目文本语义、分层语义标识、认知增强语义、题目难度、认知层级、历史作答结果和位置状态共同组成。该表示为后续目标条件认知证据检索提供统一输入，使模型能够根据当前目标题目从学生历史交互中选择具有解释价值的认知证据。

\noindent\textbf{面向跨概念迁移的协同与图结构先验}

除题目文本和知识属性外，学生群体作答行为中还包含大量题目关联信息。若多个学生在若干题目上的作答表现呈现相似模式，则这些题目之间可能存在隐含的知识联系、难度联系或学习路径联系。本文据此构建协同先验，用于捕捉题目在群体行为层面的关联。

协同先验不直接来自题目文本，而是来自学生—题目交互数据所反映的共现关系和作答模式。对于题目 \(q_j\)，协同表示 \(a_j\) 描述其在学生群体学习行为中的相对位置。该表示补充了文本语义和知识点标签难以覆盖的行为关联，使检索模块能够发现与目标题目在群体表现上相关的历史题目。

协同先验的作用在于扩展证据来源。某些题目在文本上并不相似，也不一定共享完全相同的知识点标签，但学生在这些题目上的表现可能高度相关。例如，一个题目可能考查某个概念的基础识别，另一个题目考查该概念在复杂情境中的应用。二者在题干表述上差异较大，但学生群体在两题上的表现关系可能揭示出前置掌握与迁移应用之间的联系。协同先验使后续检索能够利用这种由真实作答行为形成的题目关联，从而为目标题目预测提供不同于文本相似和知识点重叠的补充证据。

本文还构建图结构与概念关系先验，用于描述题目和知识概念之间的结构性联系。在知识追踪任务中，题目通常关联一个或多个知识概念，而概念之间可能存在共现、前置、包含或相邻关系。仅依赖直接知识点重叠会遗漏许多间接相关题目。

为此，本文将题目—概念关系和概念间结构关系组织为图结构信息，并由此得到题目的结构表示 \(g_j\)。该表示用于补充直接知识点标签，使检索模块能够识别与目标题目存在间接概念联系或结构邻近关系的历史交互。

综合上述先验，本文将题目侧信息表示为结构化认知先验：
\[
r_j=\left(\mathbf e_j^{\mathrm{sem}},\mathbf a_j,\mathcal N_j^{\mathrm{graph}},\mathcal N_j^{\mathrm{collab}},\mu_j,\ell_j\right).
\]

其中，\(\mathbf e_j^{\mathrm{sem}}\) 表示由题干文本、层级语义标识及认知属性共同形成的题目语义表示。\(\mathbf a_j\) 表示题目在学生群体作答行为中的关联表示。\(\mathcal N_j^{\mathrm{graph}}\) 与 \(\mathcal N_j^{\mathrm{collab}}\) 分别表示知识结构与协同行为邻域。\(\mu_j\) 和 \(\ell_j\) 分别表示题目难度与认知层级。

该结构化先验不直接替代学生历史序列，而是为后续目标条件认知证据检索提供题目侧认知基础。



通过将题目语义、知识结构、群体行为和认知属性统一到先验建模中，本文为每个候选历史交互赋予可检验的迁移描述。后续模块因此能够围绕当前未见目标题目进行条件化证据选择和上下文组织。

未见概念归纳评估协议要求，任何依赖作答记录的难度、协同或图统计均只由 \(\mathcal{C}_{\mathrm{train}}\) 的交互构建。测试概念的标签及其交互不得参与该类统计的估计。因而，未见目标题目可使用其作答前的文本、标注和外生结构属性定位迁移路径，但不能通过目标概念的行为统计获得隐性监督。

\noindent\textbf{面向未见目标的两阶段迁移证据检索}



对于关联未见概念的目标题目，学生历史中不存在可由该目标概念 ID 直接索引的拟合监督。因而，检索模块的职责不是从长序列中机械压缩信息。它需要以目标题目先验为条件，识别学生在既有概念上的哪些交互能够作为该目标的迁移证据。

给定学生 \(u\)、预测时刻 \(t\) 以及目标题目 \(q_{u,t}\)，模块仅从目标前历史 \(H_{u,<t}\) 中选择与 \(q_{u,t}\) 存在可解释关系的证据集合 \(E_{u,t}\)。该集合随后输入上下文构造与预测模块。这样，目标概念的语义、认知或结构描述决定\emph{检索什么}，而学生在已见概念上的作答结果决定这些证据对该目标提供何种支持。
设学生 \(u\) 的目标前历史为
\[
H_{u,<t}=\{(q_1,y_1),(q_2,y_2),\ldots,(q_{t-1},y_{t-1})\}
\]

其中 \(q_i\) 表示第 \(i\) 次历史交互中的题目，\(y_i\in\{0,1\}\) 表示学生在该题上的作答结果。两阶段认知证据检索模块以 \((u,t,q_{u,t},H_{u,<t})\) 为输入，输出证据集合
\[
E_{u,t}=\{e_{u,t}^{1},e_{u,t}^{2},\ldots,e_{u,t}^{K}\}
\]

其中每个证据 \(e_{u,t}^{k}\) 对应一个历史交互及其与目标题目之间的认知关系。与一般序列截断或最近交互选择不同，本文的检索过程由目标题目驱动。同一段学生历史面对不同目标题目时会产生不同的证据集合，因为历史题目的证据价值取决于其与当前目标题目之间的知识联系、语义联系、认知角色以及学生在该题上的作答状态。

目标条件性是本文证据检索的核心约束。历史交互并非孤立地反映学生能力，而是在特定目标题目下体现不同的认知意义。

例如，学生曾经答对某一基础题目，在预测同一知识链条上的进阶题时提供前置支撑。而在预测另一个弱相关题目时，该交互的证据价值明显降低。相应地，学生曾经答错某一语义相近或结构相关题目，在预测当前题时能够提示潜在风险。本文通过两阶段检索机制，将这种历史交互—目标题目之间的条件关系显式建模，使后续预测不依赖完整历史拼接，而依赖经过筛选和组织的认知证据。

\noindent\textbf{候选召回}

第一阶段以目标题目前的有限历史窗口构成候选集合：
\[
C_{u,t}=\{(q_{u,i},y_{u,i})\mid t-W\leq i<t\}.
\]

对于候选集合中的每条历史交互，模型不把知识概念、语义、认知层级、协同行为和图结构视为相互独立的召回通道，而是将其作为刻画历史题目与目标题目关系的多源信号。第一阶段据此计算目标相关性分数，并保留高分候选，为后续证据选择提供覆盖不同认知关系的候选空间。

\noindent\textbf{证据精排}

第二阶段在候选集合 \(C_{u,t}\) 上评估证据价值。对于候选证据 \(e_i=(q_i,y_i)\)，其相对于目标题目 \(q_{u,t}\) 的支持分数定义为
\[
s_{i,t}=\lambda_{\mathrm{k}}s^{\mathrm{k}}_{i,t}+\lambda_{\mathrm{s}}s^{\mathrm{s}}_{i,t}+\lambda_{\mathrm{g}}s^{\mathrm{g}}_{i,t}+\lambda_{\mathrm{c}}s^{\mathrm{c}}_{i,t}+\lambda_{\mathrm{y}}s^{\mathrm{y}}_{i,t}+\lambda_{\mathrm{r}}s^{\mathrm{r}}_{i,t}.
\]



其中，\(s^{\mathrm{k}}_{i,t}\) 表示概念关联，\(s^{\mathrm{s}}_{i,t}\) 表示语义与认知层级关联，\(s^{\mathrm{g}}_{i,t}\) 表示知识结构关联，\(s^{\mathrm{c}}_{i,t}\) 表示协同行为关联。\(s^{\mathrm{y}}_{i,t}\) 表示历史作答结果对当前预测的支持方向，\(s^{\mathrm{r}}_{i,t}\) 表示时间衰减项。

各项系数在拟合集内确定或按统一配置预设，并在测试阶段固定。对未见概念目标而言，前四类关系项由目标题目的作答前先验与已见历史题目的先验计算，而 \(s^{\mathrm{y}}_{i,t}\) 仅编码学生在历史题目上的已观测结果。评分过程从不使用目标概念的作答标签。该机制将题目关系、学生表现、认知角色和历史位置统一为面向跨概念迁移的排序依据。

前置支撑表示历史题目与目标题目之间存在基础到进阶的关系，学生在前置题上的表现用于反映完成目标题目所需基础能力的掌握情况。同质迁移表示历史题目与目标题目在知识概念、题型结构或认知要求上高度接近，学生在此类题目上的表现直接反映其对相似任务的处理能力。

语义近邻表示历史题目与目标题目在文本语义或问题表述上相近，该类证据用于补充显式知识点标注无法覆盖的题目相似性。高阶迁移表示历史题目与目标题目之间存在跨概念或跨层级的关联，反映学生能否将已有知识迁移到更复杂或不同表述的任务中。

图结构补充表示历史题目与目标题目通过概念图或结构关系相连，该类证据提供直接知识点重叠之外的结构性联系。协同补充表示两类题目在学生群体作答行为中表现出相关性，该类证据用于反映由真实学习行为形成的隐含题目关联。

在计算候选证据价值时，本文同时考虑作答状态的方向性。对于与目标题目关系紧密的历史题目，正确作答通常提供正向支持，错误作答则提示潜在风险。对于较高难度或高阶迁移关系下的历史题目，作答结果还反映学生在复杂认知要求下的稳定性。

本文将这种方向性纳入支持分数，使精排结果不仅区分相关和不相关，也区分历史证据对当前预测的支持方向。由此，证据检索过程不等同于题目相似度排序，而是面向知识追踪预测的认知证据选择。

\noindent\textbf{覆盖与冗余控制}

仅按照支持分数选择排名靠前的候选交互，容易得到大量来源相近、角色重复或语义高度重合的历史题目。这样的证据集合虽然局部相关性较高，但不能充分覆盖学生当前预测所需的多方面认知信息。

为此，本文在最终证据选择中引入覆盖与冗余控制，使输出集合既保留高支持强度证据，又覆盖不同认知角色和关系来源。覆盖控制要求最终证据集合包含多种认知关系来源。对于一个目标题目，模型不仅关注同知识点题目，也关注语义近邻、前置支撑、图结构邻近和协同行为相关题目。

不同角色的证据从不同角度描述学生状态：同质迁移强调相似任务表现，前置支撑强调基础能力，语义近邻强调题干表达相关性，高阶迁移强调跨层级能力，图结构补充强调概念结构联系，协同补充强调群体行为关联。

通过保留多种角色的代表性证据，最终证据集合获得更完整的目标条件覆盖。冗余控制要求最终证据集合避免被少数相似题目占据。

若多个候选历史题目具有高度相近的语义表示、相同的知识概念和相同的证据角色，模块优先保留支持分数更高、历史位置更有代表性或作答状态更具解释性的交互。这样可以降低重复证据对上下文的占用。该机制使最终证据集合在有限长度内保留更多互补信息，为后续分层认知上下文构造提供稳定输入。

形式化地，本文以贪心方式构造最终证据集合：
\[
E_{u,t}=\arg\max_{|E|=K}\left[\sum_{e_i\in E}s_{i,t}+\lambda_{\mathrm{cov}}\operatorname{Cov}(E,q_{u,t})-\lambda_{\mathrm{red}}\operatorname{Red}(E)\right],
\]



其中，\(\operatorname{Cov}(E,q_{u,t})\) 衡量证据集合对不同迁移关系来源和目标能力要求的覆盖程度。\(\operatorname{Red}(E)\) 衡量证据之间的语义、概念和角色重复程度。该约束防止模型仅凭单一表面相似路径对未见目标作出判断，使集合在有限容量下同时保留前置支撑、同质迁移和结构补充等互补证据。历史窗口长度 \(W\)、候选规模和最终证据数 \(K\) 在实验设置中报告。

最终，目标条件认知证据检索模块输出证据集合 \(E_{u,t}\)。每条证据包含历史题目、历史作答结果、证据角色、支持分数、历史位置以及其与目标题目之间的关系来源。

该输出并不是简单的历史题目列表，而是经过目标条件召回、证据价值精排和覆盖冗余控制后形成的结构化认知证据集合。后续分层认知上下文构造模块基于该集合组织学生侧认知状态与局部证据块，并将其传递给基于冻结 LLM 的 soft-slot 概率预测模块，用于计算学生在目标题目上的作答正确概率。

\subsection{分层认知上下文构造}

\noindent\textbf{从跨概念离散证据到结构化迁移上下文}



两阶段检索输出由既有概念历史提供的离散迁移证据。若只按时间或检索分数拼接这些交互，预测器仍需自行区分一般学生状态、当前目标支持证据及其跨概念关系，局部证据也容易堆叠而缺少迁移结构。为此，本文构建分层认知上下文，将检索结果组织为包含全局认知状态、局部迁移证据块和目标题目表示的结构化表示。



给定学生 \(u\)、预测时刻 \(t\)、目标题目 \(q_{u,t}\) 及检索得到的证据集合 \(E_{u,t}\)，该模块生成学生侧认知上下文表示 \(x^{ctx}_{u,t}\) 与目标题目侧表示 \(x^{tar}_{t}\)。前者描述学生在既有概念历史中呈现的表现、证据结构和状态特征。后者描述当前未见题目的语义、层级、认知及可用关系先验。

二者不能相互替代。学生侧表示提供学生已表现出的能力证据，目标题目侧表示规定这些证据应被迁移到何种能力要求上。



后续冻结 LLM 的 soft-slot 概率预测模块分别接收这两类表示，并通过投影器将其注入 context soft slots 和 target soft slots。本文构造的上下文既不是完整历史的自然语言展开，也不是检索条目的简单列表。它显式保留被选中交互、交互与未见目标之间的关系，以及学生在这些交互上的作答表现，从而为概率预测提供有来源的迁移条件。

分层认知上下文因此承担连接检索模块与预测模块的作用：检索模块负责从目标前历史中选择相关证据，上下文构造模块负责将这些证据转化为具有层次结构和数值特征的认知状态表示，预测模块据此计算当前目标题目的作答正确概率。

\noindent\textbf{目标条件全局认知状态}

全局认知状态用于概括学生在当前目标题目条件下的整体学习状态。与单条历史证据相比，全局状态强调证据集合所反映的综合趋势，包括近期作答表现、风险水平、动态状态以及主要证据角色分布。本文将检索证据集合 \(E_{u,t}\) 与学生历史交互中的状态信息结合，生成全局认知状态 \(G_{u,t}\)。该状态不替代局部证据，而是为局部证据提供整体背景。

具体地，本文仅对检索得到的目标前证据按照目标相关性加权聚合：
\[
{G_{u,t}=\sum_{k=1}^{K}\alpha_{k,t}\left[\mathbf e_{i_k}^{\mathrm{sem}};y_{u,i_k};\boldsymbol{\rho}_{i_k,t};\delta_{i_k,t}\right],\qquad \alpha_{k,t}=\frac{\exp(s_{i_k,t})}{\sum_{j=1}^{K}\exp(s_{i_j,t})},}
\]

其中，\(i_k\) 表示第 \(k\) 条检索证据在目标前历史中的原始位置，\(\boldsymbol{\rho}_{i_k,t}\) 表示该历史交互相对于目标题目的认知关系角色，\(\delta_{i_k,t}\) 表示其历史位置信息。由此，\(G_{u,t}\) 用于描述整体学习趋势与风险，局部证据块则保留单条交互对当前预测的具体支持来源。

首先，全局认知状态包含近期趋势信息。学生知识状态具有时间演化特征，近期交互通常更直接地反映当前掌握情况。本文在构造全局状态时纳入近期作答正确率、近期错误模式以及目标前历史中的表现变化，使上下文能够描述学生在预测时刻附近的状态走向。该信息用于区分长期历史中的偶然表现与当前阶段的稳定趋势。

其次，全局认知状态包含风险水平信息。风险水平由学生在相关证据上的作答结果、题目难度以及证据与目标题目的支持关系共同决定。当学生在与目标题目高度相关的历史题目上出现错误，或者在前置支撑题目上表现不稳定时，全局状态会反映较高的预测风险。当学生在多类相关证据上保持正确表现时，全局状态则体现较强的正向支持。该风险描述为后续预测提供目标条件下的整体判断基础。

第三，全局认知状态包含动态学生状态与主要证据角色。动态状态描述学生在目标前历史中的局部掌握变化，主要证据角色描述当前证据集合由哪些认知关系主导。例如，当证据集合主要由同质迁移和语义近邻构成时，上下文强调学生在相似题目上的表现。当证据集合主要由前置支撑和高阶迁移构成时，上下文强调基础能力到复杂任务之间的迁移关系。

通过记录证据角色分布，全局状态能够为预测模块提供关于证据来源结构的概括信息。

\noindent\textbf{局部迁移证据块}

局部证据块用于保留每条检索证据的细粒度信息。对于证据集合 \(E_{u,t}\) 中的每条证据 \(e_{u,t}^{k}\)，本文构造对应的局部证据块 \(B_{u,t}^{k}\)。每个证据块包含历史题目表示、历史作答结果、证据关系、支持强度、历史位置以及与目标题目之间的角色标注。与全局认知状态不同，局部证据块不对证据集合进行整体压缩，而是保留单条证据对当前预测的具体贡献来源。

局部证据块首先记录历史题目及其题目侧认知先验。历史题目的文本语义、分层语义标识、认知增强语义、协同表示和图结构关系共同描述该题在多源认知空间中的位置。其次，证据块记录学生在该历史题目上的作答结果。作答标签决定该证据对当前预测的支持方向：对于与目标题目关系紧密的历史题目，正确作答提供正向证据，错误作答提示潜在风险。

第三，证据块记录该历史交互与目标题目之间的证据关系，包括前置支撑、同质迁移、语义近邻、高阶迁移、图结构补充和协同补充等角色。该角色信息使预测模块能够区分不同证据的认知含义，而不是将所有历史题目视为同质输入。

局部证据块还包含支持强度和历史位置信息。支持强度来自证据精排阶段，反映该历史交互相对于目标题目的综合证据价值。历史位置信息描述该交互在学生目标前序列中的时间位置，用于区分近期证据和较早证据。对于知识追踪任务，时间位置并不简单等同于证据价值，但它为理解学生状态变化提供必要背景。因此，本文在证据块中同时保留关系强度与历史位置，使局部证据既具有目标相关性，也具有序列状态含义。
通过局部证据块组织，检索结果被转化为结构化证据序列：
\[
B_{u,t}=\{B_{u,t}^{1},B_{u,t}^{2},\ldots,B_{u,t}^{K}\}
\]

该序列中的每个元素均对应一条目标条件证据，并包含其题目属性、作答状态、证据角色和支持信息。后续上下文表示在保留局部证据差异的同时，与全局认知状态共同描述学生当前的知识掌握情况。

\noindent\textbf{结构化认知摘要}

在全局状态和局部证据块之外，本文进一步引入结构化认知摘要，对检索证据进行目标条件下的紧凑压缩。该步骤使用冻结 LLM 读取目标题目、学生近期统计摘要和已筛选的目标条件证据，但其输出被限制为受约束摘要字段，而不是自由文本预测。具体而言，摘要包含稳定掌握点、持续薄弱点、迁移状态、风险等级、证据质量和简短证据说明，用于描述检索证据在当前目标题目上的条件含义。

形式化地，结构化认知摘要记为
\[
D^{\mathrm{sum}}_{u,t}=\{\mathcal M_{u,t},\mathcal W_{u,t},\tau_{u,t},r_{u,t},\eta_{u,t},d_{u,t}\},
\]
其中 \(\mathcal M_{u,t}\) 和 \(\mathcal W_{u,t}\) 分别表示由证据支持的掌握点与薄弱点，\(\tau_{u,t}\) 表示迁移状态，\(r_{u,t}\) 表示风险等级，\(\eta_{u,t}\) 表示证据质量，\(d_{u,t}\) 表示简短证据说明。

该摘要不作为最终标签，也不参与开放式解释生成。其作用是把多条检索证据压缩为可编码的证据压缩信号。随后，摘要说明被编码为上下文向量，摘要中的结构化字段被转换为显式特征，并与全局状态、局部证据块共同构成学生侧上下文输入。

为保证该步骤可复现，本文将结构化认知摘要实现为受 schema 约束的证据压缩器，而非开放式文本生成模块。摘要生成器的输入仅包含目标题目的可观察信息、目标概念集合、学生目标前近期模板统计摘要以及已筛选证据 \(E_{u,t}\) 的证据编号、角色、支持强度、历史作答结果和题目文本。输入不包含当前目标作答标签。

输出被限制为一个 JSON object。该对象必须包含六个顶层字段：\texttt{mastered\_concepts}、\texttt{weak\_concepts}、\texttt{transfer\_state}、\texttt{risk\_level}、\texttt{evidence\_quality} 和 \texttt{diagnosis}。

其中，掌握点和薄弱点字段最多各包含 3 个条目。每个条目由概念名、证据编号集合和置信等级组成。风险等级、证据质量和置信等级均使用离散等级 \{低，中，高\}。

实现中通过 JSON-only 系统提示、显式用户提示、\texttt{response\_format=json\_object}、温度 0.0 和最大输出长度 256 tokens 约束解码。解析失败或字段不合法时最多触发三次重试。仍不合法的摘要不会进入后续编码。

在编码阶段，合法摘要首先被转写为稳定掌握、薄弱风险、诊断说明和辅助状态等短文本通道，并与目标条件上下文一起形成摘要侧嵌入。同时，摘要字段被转换为固定维度结构特征，包括掌握点数量、薄弱点数量、是否存在掌握点、是否存在薄弱点、诊断长度归一化值、风险等级 one-hot 和证据质量 one-hot。

上述摘要嵌入和摘要派生结构特征作为学生侧上下文特征的一部分进入 context projector。因此，结构化认知摘要在本文中是可验证的中间表示，而不是需要后验解析的自由文本解释。

\noindent\textbf{面向冻结预测器的上下文表示}

分层认知上下文最终被转换为面向 soft-slot 概率预测模块的结构化特征。本文将上下文表示分为学生侧认知上下文表示和目标题目侧表示两部分。学生侧认知上下文表示由全局认知状态、局部证据块、结构化认知摘要、证据结构特征和数值状态特征共同构成。目标题目侧表示由题目文本语义、分层语义标识、认知增强语义、协同向量以及相关认知属性构成。两类表示分别对应学生当前状态与当前预测目标，使预测模块能够在同一任务约束下同时获得学生具备什么状态和当前题目要求什么能力的信息。

学生侧认知上下文嵌入用于整合全局状态和局部证据块，并吸收结构化认知摘要中的压缩信息。该嵌入吸收检索证据中关于历史题目、作答结果、证据角色和支持强度的信息，形成目标条件下的学生状态表示。

证据结构特征用于保留证据集合的组成方式，例如主要证据角色、证据数量、支持分数分布以及不同关系来源的覆盖情况。摘要派生特征则保留掌握点数量、薄弱点数量、迁移状态、风险等级和证据质量等压缩后的证据信号。数值状态特征用于描述近期表现、风险水平、动态趋势等统计状态。上述信息共同构成学生侧认知上下文特征，使后续预测不仅依赖语义嵌入，也利用显式结构和数值状态。

目标题目侧表示用于描述当前预测对象。本文将目标题目的文本语义表示用于刻画题目内容，将分层语义标识用于刻画其在语义层级中的位置，将认知增强语义表示用于刻画题目难度、认知方向或结构属性，将协同向量用于刻画该题在学生群体作答行为中的关联位置。目标题目侧表示与学生侧上下文表示相互配合：前者描述预测目标，后者描述学生在该目标条件下的证据状态。两者共同决定 soft-slot 概率预测模块接收的非文本结构化输入。
分层认知上下文构造模块输出
\[
X_{u,t}=(x^{ctx}_{u,t},x^{tar}_{t})
\]

其中 \(x^{ctx}_{u,t}\) 表示学生侧认知上下文特征，\(x^{tar}_{t}\) 表示目标题目侧多源先验特征。该输出将离散检索证据压缩和组织为适合预测模块处理的结构化表示，并通过结构化认知摘要补充目标条件证据压缩信息，从而避免将完整历史交互直接堆叠到提示文本中。



后续冻结 LLM 的 soft-slot 概率预测模块将 \(x^{ctx}_{u,t}\) 与 \(x^{tar}_{t}\) 分别映射为 context soft slots 和 target soft slots，再在标签约束下计算学生对目标题目作答正确的概率。

由此，分层认知上下文构造承担承上启下的作用。它把已见概念上的学生表现、跨概念证据角色和未见目标的可观察先验整合为统一表示，同时保留每一类信息在未见概念推断中的职责边界。

\subsection{基于冻结 LLM 的 soft-slot 概率预测}

\noindent\textbf{冻结 LLM 的受约束表示空间}



经过分层认知上下文构造和结构化认知摘要生成后，模型获得学生侧认知上下文特征 \(x^{ctx}_{u,t}\) 和目标题目侧多源先验特征 \(x^{tar}_{t}\)。前者概括学生在既有概念上的、与当前目标相关的证据状态及摘要压缩后的证据信号。后者提供未见目标的语义、层级、认知与关系描述。

基于这两类输入，本文构建冻结 LLM 的 soft-slot 概率预测模块。该模块学习在目标条件下将既有表现映射为对未见概念题目的正确率估计 \(\hat{p}_{u,t}\)。



本文将冻结 LLM 作为受约束的条件表示与标签判别空间，而非作为自由生成器或外部知识来源。LLM 主体参数保持不训练。训练仅更新学生侧上下文投影器和目标题目侧投影器，使前序模块输出的结构化认知证据能够进入 LLM 的连续输入空间。

该约束具有两层作用。首先，它限制模型通过改变语言模型参数记忆测试概念的监督信息。其次，它使未见概念上的预测改进主要由目标先验、证据检索和条件表示之间的协同来解释，而不是由无边界文本生成来解释。

设冻结 LLM 的参数为 \(\Theta\)，训练过程中保持 \(\Theta\) 不变。学生侧上下文投影器记为 \(P_{\phi}^{ctx}\)，目标题目侧投影器记为 \(P_{\psi}^{tar}\)，其中 \(\phi\) 和 \(\psi\) 为需要学习的参数。对于每个预测样本 \((u,t,q_{u,t})\)，投影器分别将学生侧认知上下文和目标题目侧先验映射为连续向量序列：
\begin{align}
Z^{ctx}_{u,t} &= P_{\phi}^{ctx}(x^{ctx}_{u,t}), \\
Z^{tar}_{t} &= P_{\psi}^{tar}(x^{tar}_{t}).
\end{align}

其中，\(Z^{ctx}_{u,t}\) 表示 context soft slots，\(Z^{tar}_{t}\) 表示 target soft slots。两类 soft slots 分别承担不同功能。context soft slots 注入学生在当前目标条件下的认知状态，target soft slots 注入当前题目的多源认知先验。通过这种设计，预测模块在不展开完整历史文本的情况下接收结构化证据，并在冻结 LLM 的表示空间中完成二分类预测。

\noindent\textbf{目标条件 soft slots 注入}



本文将 soft slots 设计为连接结构化认知特征与冻结 LLM 的中间表示。不同于将历史交互、检索证据和统计特征转写为长文本并要求 LLM 自行推断，soft slots 将学生在已见概念上的迁移证据和未见目标的先验分别编码为连续条件输入。该设计保留检索与上下文构造输出的结构，同时避免完整历史拼接造成的长度增长、证据混淆及自由文本解析误差。

context soft slots 来源于学生侧认知上下文特征 \(x^{ctx}_{u,t}\)。该部分信息包含学生历史交互中被检索出的证据、证据角色、支持强度、近期状态和风险水平，也包含由结构化认知摘要派生的掌握点、薄弱点、迁移状态和证据质量信息。学生侧投影器 \(P_{\phi}^{ctx}\) 将这些异构特征映射为固定长度的连续向量序列，使其能够进入 LLM 的嵌入层输入空间。

context soft slots 的作用不是复述学生历史，而是向预测模块提供已经筛选、压缩和组织后的认知状态条件。对于同一学生，不同目标题目会触发不同的证据检索结果。因此，对应的 context soft slots 也随目标题目变化，体现知识追踪预测的目标条件性。

target soft slots 来源于目标题目侧多源先验特征 \(x^{tar}_{t}\)。该部分信息包含题目文本语义、分层语义标识、认知增强语义、协同关系和题目认知属性。目标题目侧投影器 \(P_{\psi}^{tar}\) 将这些题目先验映射为连续向量序列，用于向冻结 LLM 注入当前题目的能力要求和题目侧背景。

target soft slots 与 context soft slots 共同构成预测条件。前者描述当前题目要求什么，后者描述学生在该目标条件下表现出什么状态。


在输入组织上，本文同时保留文本提示模板和连续 soft slots。提示模板只声明知识追踪任务与候选标签约束，不注入目标概念的答案或依赖未验证的外部事实。soft slots 承载实际用于预测的结构化认知证据。提示模板不承担完整历史拼接功能，也不要求 LLM 从冗长自然语言历史中自行筛选证据。

输入到冻结 LLM 的条件信息由三部分组成：任务提示文本、context soft slots 和 target soft slots，其整体输入记为
\[
I_{u,t}=[T;\,Z^{ctx}_{u,t};\,Z^{tar}_{t}]
\]

其中 \(T\) 表示任务说明与标签约束提示，\(Z^{ctx}_{u,t}\) 和 \(Z^{tar}_{t}\) 分别表示学生侧与题目侧连续条件输入。该输入结构将知识追踪任务约束、学生认知状态和目标题目先验同时注入预测过程。

\noindent\textbf{受标签约束的 logits 概率预测}



本文将未见概念上的知识追踪建模为受信息边界约束的二分类概率估计，而不是自然语言生成任务。对于输入 \(I_{u,t}\)，冻结 LLM 前向计算后在候选标签位置产生词表 logits。预测模块仅读取正确与错误标签的 logits，并通过 softmax 得到学生作答正确的概率。因而，LLM 的作用是对已编码的学生证据与目标先验进行条件判别，输出中不存在需要解析的自由文本。



其中，\(p_{u,t}=P(y_{u,t}=1\mid H_{u,<t},q_{u,t},a(q_{u,t}))\) 表示在目标前历史与目标题目可观察先验条件下的理论正确概率，\(\hat p_{u,t}\) 表示模型对该概率的估计。设 \(o^{+}_{u,t}\) 和 \(o^{-}_{u,t}\) 分别表示正确与错误标签对应的 logit，则
\[
\hat p_{u,t}=f_{\Theta,\phi,\psi}(I_{u,t})=\frac{\exp(o^{+}_{u,t})}{\exp(o^{+}_{u,t})+\exp(o^{-}_{u,t})}.
\]

相应地，作答错误的概率为：
\[
P(y_{u,t}=0 \mid H_{u,<t}, q_{u,t}) = \frac{\exp(o^{-}_{u,t})}{\exp(o^{+}_{u,t}) + \exp(o^{-}_{u,t})}
\]

这种输出方式使模型预测结果直接对应知识追踪评价所需的概率分数。与要求 LLM 生成自然语言答案或结构化 JSON 再进行解析不同，本文不依赖生成文本的后处理过程。模型直接使用标签 logits 完成二分类概率计算。预测结果 \(\hat{p}_{u,t}\) 既用于训练阶段的损失计算，也用于测试阶段的 AUC、ACC、F1 和 RMSE 等指标评估。

训练阶段，给定样本真实标签 \(y_{u,t} \in \{0,1\}\)，本文采用二分类交叉熵损失优化投影模块参数：

\[
\mathcal{L} = -\sum_{(u,t)} \left[
    y_{u,t} \log \hat{p}_{u,t}
    + (1-y_{u,t}) \log(1-\hat{p}_{u,t})
\right]
\]



在该目标下，冻结 LLM 参数 \(\Theta\) 不更新，优化对象为学生侧上下文投影器参数 \(\phi\) 和目标题目侧投影器参数 \(\psi\)。训练促使投影器学习如何将已见概念上的学生侧认知上下文与未见目标的多源先验映射到适合标签判别的连续输入空间。

由于 LLM 主体保持冻结，任务适配发生在证据到 soft slots 的映射层。这使模型必须依靠可迁移的条件表示，而不能通过更新 LLM 参数为测试概念建立专有记忆。

通过 soft slots 注入，模型将检索到的目标相关认知证据转化为连续条件输入，并在冻结 LLM 的标签空间中完成知识追踪预测。该连续条件输入设置与冻结语言模型中的 prompt tuning 和 prefix tuning 同属参数高效连续提示思想。

本文的 soft slots 由认知证据特征投影得到，并服务于二分类标签 logits 预测\citep{lester2021power,li2021prefixtuning}。该过程保持了知识追踪任务的时序约束和概率输出形式，同时避免将学生完整历史直接交给 LLM 进行无约束生成。

\setcounter{section}{3}
\section{实验}
\subsection{实验设置}


\subsubsection{数据集与划分}

本文在三个知识追踪数据集上评估所提出的 Cognitive-RAG KT 方法，分别为 MOOCRadar~\citep{yu2023moocradar}、XES3G5M~\citep{liu2023xes3g5m} 和 FoundationalASSIST~\citep{worden2026foundationalassist}。三个数据集均被统一整理为学生--题目交互序列形式，每条交互记录包含学生标识、题目标识、作答结果以及题目关联的知识概念。模型的预测目标是在给定学生历史交互与目标题目认知先验的条件下，估计学生在目标题目上作答正确的概率。

为保证实验设置与方法目标一致，本文采用面向目标题目的预测样本构造方式。对于学生 \(u\) 在时间步 \(t\) 的目标题目 \(q_{u,t}\)，其历史输入仅由 \(H_{u,<t}\) 构成，真实标签为 \(y_{u,t} \in \{0,1\}\)，其中 \(y_{u,t}=1\) 表示学生作答正确，\(y_{u,t}=0\) 表示作答错误。目标题目自身的文本、概念、语义向量、协同向量等信息用于构造题目侧认知先验，但目标题目的真实作答结果不作为 prompt 或 embedding 输入。

\begin{table}[H]
\centering
\small
\caption{数据集总体统计}
\label{tab:dataset-split-statistics}
\resizebox{\textwidth}{!}{%
\begin{tabular}{lrrrr}
\toprule
\textbf{数据集} & \textbf{学生数} & \textbf{题目数} & \textbf{概念数} & \textbf{交互数}  \\
\midrule
MOOCRadar & 14,224 & 9,383 & 9,985 & 898,933  \\
XES3G5M & 4,833 & 3,903 & 575 & 609,867  \\
FoundationalASSIST & 5,000 & 2,492 & 143 & 999,467  \\
\bottomrule
\end{tabular}%
}
\end{table}

\subsubsection{未见概念归纳评估协议}

知识追踪中的冷启动问题并不只表现为新学生或新题目，也可能表现为训练阶段尚未出现的知识概念。本文将其作为未见概念归纳评估场景来处理，而不将其重新定义为独立的新任务：测试目标至少关联一个训练监督中未出现的知识概念，模型因而不能仅依赖对既有概念或题目 ID 的经验进行预测。该设定考察模型能否结合学生既往学习历史与目标题目的可获得先验，将已学习到的认知表征迁移至未见知识概念。

具体而言，本文以题目关联的知识概念为划分依据构造训练集和测试集。对于测试目标 \(q_{u,t}\)，其关联概念集合至少包含一个训练阶段未见的测试概念，即 \(C(q_{u,t})\cap\mathcal{C}_{\mathrm{test}}\neq\emptyset\)。若一道题同时关联已见概念和未见概念，只要其包含任一测试概念，即纳入测试目标集合；因此该协议考察模型面对未见能力要求时的归纳预测能力。

测试阶段仅使用目标题目发生前的学生历史交互，以及作答前可获得的题目文本、概念信息和静态属性。依赖作答记录构建的难度、协同关系和图结构统计均由训练交互生成，并在测试阶段固定。

\begin{table}[H]
\centering
\small
\caption{未见概念归纳评估的划分统计}
\label{tab:cold-start-split-statistics}
\resizebox{\textwidth}{!}{%
\begin{tabular}{lrrrrr}
\toprule
\textbf{数据集} & \textbf{训练交互} & \textbf{测试交互} & \textbf{训练概念} & \textbf{测试概念} & \textbf{未见测试概念} \\
\midrule
FoundationalASSIST & 328,433 & 666,034 & 27 & 127 & 113 (89.0\%) \\
XES3G5M & 108,319 & 496,715 & 99 & 524 & 476 (90.8\%) \\
MOOCRadar & 66,512 & 818,197 & 138 & 2,703 & 2,642 (97.7\%) \\
\bottomrule
\end{tabular}%
}
\end{table}

表~\ref{tab:cold-start-split-statistics} 显示，三个数据集的测试目标均全部包含至少一个训练阶段未见概念，且测试概念中未见概念占比较高。为保持评估边界一致，依赖作答记录构建的题目难度、协同关系和图结构统计均仅由训练交互估计；测试标签、测试交互及其派生统计不参与题目先验构建或模型选择。

\subsubsection{对比模型}

本文主实验比较已有知识追踪基线模型与本文方法，结果仅报告在相同数据划分和相同测试样本上完成评估的模型。外部对比模型包括 DKT~\citep{piech2015deep}、SAKT~\citep{pandey2019sakt}、SAINT~\citep{choi2020saint}、AKT~\citep{ghosh2020akt}、DIMKT~\citep{shen2022dimkt} 和 QIKT~\citep{chen2023qikt} 等代表性模型。

此外，本文还比较 TCKT~\citep{huang2024tckt}、SimpleKT~\citep{liu2023simplekt}、SparseKT~\citep{huang2023sparsekt}、DenoiseKT~\citep{deng2025denoisekt}、RobustKT~\citep{guo2025robustkt} 和 KeenKT~\citep{li2026keenkt}。各模型的类别与输入信息见附录表~\ref{tab:appendix-comparison-models}。

传统 KT 对比模型主要基于学生历史交互序列进行预测，不使用 Cognitive-RAG 检索得到的认知证据。Cognitive-RAG KT 使用题目侧多源认知先验和学生侧认知上下文特征；后者包含检索证据、结构化证据属性和结构化认知摘要派生特征。模型将这些多源信息投影为 soft slots，输入冻结的 Qwen3-8B\citep{yang2025qwen3,qwen2025qwen38b}，再读取标签 logits 并通过二分类 softmax 得到 \(\hat{p}_{u,t}\)，而不是通过自由文本生成完成预测。

因此，本文比较的重点不是所有模型在输入信息上的完全同构，而是在相同数据划分和相同测试目标下，检验不同建模范式对未见概念归纳预测的适应能力。结果解释也相应限定为：认知检索证据与冻结 LLM 表征空间结合后，对知识追踪概率预测是否带来增益。

\subsubsection{实现细节}

本文依次执行目标条件认知证据检索、分层认知上下文构造、结构化认知摘要生成、soft-slot 特征转换和概率预测。候选证据由目标题目的知识、语义、协同和图结构关系召回，并结合相关性、覆盖性和冗余控制完成精排。结构化认知摘要将检索证据压缩为掌握点、薄弱点、迁移状态、风险等级和证据质量等字段，并作为学生侧上下文特征进入投影器。

预测模块采用冻结的 Qwen3-8B，仅训练学生侧与目标题目侧投影器。主实验同时使用 context soft slots 和 target soft slots，并通过标签 logits 的二分类 softmax 计算 \(\hat{p}_{u,t}\)\citep{yang2025qwen3,qwen2025qwen38b,qwen2025qwen3embedding,qwen2025qwen3reranker}。完整的检索设置、结构化认知摘要 schema、soft-slot 预测配置和训练配置见附录表~\ref{tab:appendix-main-method-config}。

\subsection{主实验结果与分析}

表~\ref{tab:main-results} 给出了 Cognitive-RAG KT 与各类知识追踪基线模型在三个数据集上的主实验结果。为避免与直接 prompt LLM 方法混淆，本文在正文中统一将所提方法称为 Cognitive-RAG KT。需要说明的是，本节仅分析主实验结果，不将整体性能差异直接归因于某一个子模块；组件贡献在后续消融实验中单独讨论。

\begin{landscape}
\begin{table}[p]
\centering
\caption{主实验结果对比}
\label{tab:main-results}
\scriptsize
\setlength{\tabcolsep}{4pt}
\renewcommand{\arraystretch}{1.12}
\resizebox{\linewidth}{!}{%
\begin{tabular}{lcccccccccccc}
\toprule
\multirow{2}{*}{\textbf{方法}} & \multicolumn{4}{c}{\textbf{FoundationalASSIST}} & \multicolumn{4}{c}{\textbf{XES3G5M}} & \multicolumn{4}{c}{\textbf{MOOCRadar}} \\
\cmidrule(lr){2-5}\cmidrule(lr){6-9}\cmidrule(lr){10-13}
 & \textbf{AUC} & \textbf{ACC} & \textbf{F1} & \textbf{RMSE} & \textbf{AUC} & \textbf{ACC} & \textbf{F1} & \textbf{RMSE} & \textbf{AUC} & \textbf{ACC} & \textbf{F1} & \textbf{RMSE} \\
\midrule
DKT & 0.5053 & 0.5053 & 0.5621 & 0.5085 & 0.5006 & 0.5013 & 0.6125 & 0.5100 & 0.4830 & 0.4932 & 0.6062 & 0.5161 \\
SAKT & 0.5974 & 0.5902 & 0.6667 & 0.5396 & 0.5636 & 0.7011 & 0.8127 & 0.4960 & 0.7090 & 0.7779 & 0.8666 & 0.4103 \\
SAINT & 0.6184 & 0.6226 & 0.7144 & 0.4857 & 0.5692 & 0.7831 & 0.8773 & 0.4211 & 0.6637 & 0.7024 & 0.8107 & 0.4476 \\
AKT & 0.6430 & 0.6360 & 0.7195 & 0.4855 & 0.6065 & 0.7503 & 0.8504 & 0.4235 & 0.7357 & 0.8099 & 0.8895 & 0.3716 \\
DIMKT & 0.6250 & 0.6388 & 0.7414 & 0.4839 & 0.5771 & 0.7607 & 0.8601 & 0.4228 & 0.7385 & 0.8345 & 0.9066 & 0.3518 \\
QIKT & 0.6250 & 0.6378 & 0.7377 & 0.4857 & 0.6049 & 0.7712 & 0.8670 & 0.4154 & 0.7806 & 0.8371 & 0.9081 & 0.3444 \\
TCKT & 0.6011 & 0.5967 & 0.6760 & 0.5130 & 0.6004 & 0.7635 & 0.8613 & 0.4193 & 0.7702 & 0.8220 & 0.8965 & 0.3559 \\
SimpleKT & 0.5981 & 0.5625 & 0.6158 & 0.5367 & 0.5726 & 0.7312 & 0.8363 & 0.4431 & 0.7582 & 0.8207 & 0.8953 & 0.3596 \\
SparseKT & 0.6067 & 0.6067 & 0.6918 & 0.5109 & 0.5810 & 0.7338 & 0.8381 & 0.4407 & 0.7387 & 0.8180 & 0.8941 & 0.3668 \\
DenoiseKT & 0.6557 & 0.6469 & 0.7340 & 0.4694 & 0.6384 & 0.7839 & 0.8768 & 0.4028 & 0.7817 & 0.8391 & 0.9094 & 0.3415 \\
RobustKT & 0.6252 & 0.6098 & 0.6835 & 0.4903 & 0.5745 & 0.7413 & 0.8448 & 0.4392 & 0.7448 & 0.8297 & 0.9029 & 0.3555 \\
KeenKT & 0.6246 & 0.6226 & 0.7054 & 0.4824 & 0.6232 & 0.7612 & 0.8577 & 0.4182 & 0.7758 & 0.8067 & 0.8849 & 0.3664 \\
\textbf{Cognitive-RAG KT} & \textbf{0.7805} & \textbf{0.7335} & \textbf{0.7998} & \textbf{0.4239} & \textbf{0.7542} & \textbf{0.8046} & \textbf{0.8873} & \textbf{0.3791} & \textbf{0.8404} & \textbf{0.8521} & \textbf{0.9111} & \textbf{0.3322} \\
\bottomrule
\end{tabular}%
}
\end{table}
\end{landscape}

\begin{figure}[H]
\centering
\includesvg[inkscapelatex=false,width=0.98\textwidth]{figures/fig2_paired_bootstrap_deltas_v2}
\caption{Cognitive-RAG KT 相对最强基线 DenoiseKT 的配对 bootstrap 差值。正值表示 Cognitive-RAG KT 优于 DenoiseKT；RMSE 以误差降低量表示。误差线表示 95\% 置信区间。}
\label{fig:paired-bootstrap-deltas}
\end{figure}

如表~\ref{tab:main-results} 所示，Cognitive-RAG KT 在 MOOCRadar、XES3G5M 和 FoundationalASSIST 上均取得最优的 AUC、ACC、F1 和 RMSE。主模型在三个测试集上的 AUC 分别为 0.8404、0.7542 和 0.7805；其 bootstrap 95\% 置信区间分别为 [0.8393, 0.8416]、[0.7525, 0.7559] 和 [0.7793, 0.7816]，说明结果在大规模未见概念测试样本上具有稳定的概率排序表现。

进一步地，图~\ref{fig:paired-bootstrap-deltas} 给出本文方法相对各数据集最强基线 DenoiseKT 的配对 bootstrap 差值。三个数据集上的 AUC、ACC、F1 差值和 RMSE 降低量均为正，且相应 95\% 置信区间均未跨越 0。该结果表明，在相同未见概念测试目标上，Cognitive-RAG KT 相比强序列基线不仅取得更高的点估计，也表现出统计上稳定的相对优势。

不同数据集上的收益形态仍存在差异。FoundationalASSIST 上的 AUC 提升最大，XES3G5M 上 AUC 与阈值指标均有稳定增益，MOOCRadar 上尽管正类比例较高，本文方法仍在 AUC 与 RMSE 上保持明显优势。结合三个数据集的结果，本文方法的收益不宜简单理解为输入信息量增加，而更应理解为将题目侧先验与学生目标前历史组织为面向当前预测的结构化证据后，提升了模型在未见概念目标上的条件判别能力。

\subsection{未见概念测试切片分析}

为进一步确认主结果不是由混合概念样本或少数易预测目标主导，本文对测试集进行未见概念切片分析。图~\ref{fig:unseen-slice-results} 将测试目标划分为纯未见概念样本和混合已见--未见概念样本，用于检查模型表现是否主要来自包含已见概念的测试目标。其中，纯未见样本指目标题目关联概念均不属于训练概念集合；混合样本指目标题目同时包含训练阶段已见概念和未见概念。

\begin{figure}[H]
\centering
\includesvg[inkscapelatex=false,width=0.98\textwidth]{figures/fig4_unseen_slice_composite_v2}
\caption{未见概念测试切片的性能与样本构成。左、中两部分分别比较纯未见概念样本和混合已见--未见概念样本上的 AUC 与 RMSE，其中 RMSE 越低越好；右侧给出对应测试切片的样本占比。}
\label{fig:unseen-slice-results}
\end{figure}

切片结果显示，纯未见概念样本在 FoundationalASSIST、XES3G5M 和 MOOCRadar 中分别占 96.4\%、94.5\% 和 69.8\%。在纯未见样本上，Cognitive-RAG KT 的 AUC 分别为 0.7810、0.7540 和 0.8439，与完整未见概念测试结果基本一致；在混合已见--未见样本上，模型也保持可比的排序性能。该结果表明，主模型表现并非主要由同时包含已见概念的混合样本驱动，而是在更严格的纯未见概念目标上仍保持稳定的排序表现。

除未见概念切片外，本文还将历史长度、检索证据数、语义距离、测试概念频率和图连通性作为描述性诊断维度。为避免主文过度展开，细粒度分层结果置于附录；其中的范围表示不同 bucket 的 AUC 最小--最大值，不表示置信区间，见附录图~\ref{fig:stratified-auc-ranges}。


\subsection{消融实验设计与结果分析}

消融实验聚焦当前结果中可直接对应的三项关键组件：目标条件认知证据检索、结构化认知摘要和结构化证据。完整模型保留三项组件；每个变体仅移除名称所示组件，其余输入和预测流程保持不变。

为避免组件比较受到样本选择差异影响，所有变体在同一固定评估清单上运行。该清单以三元组 \((u,t,q_{u,t})\) 对应的 \texttt{sample\_id} 对齐学生、目标时间步、目标题目和测试标签；它只用于固定评估样本与计算指标，不参与证据生成或概率预测。变体定义与共享评估协议见附录表~\ref{tab:appendix-ablation-design}。其中，结构化认知摘要指冻结 LLM 生成的受约束摘要字段及其派生特征，而不是开放式标签生成过程。

\subsubsection{关键组件消融结果}

表~\ref{tab:ablation_study} 报告目标条件认知证据检索、结构化认知摘要和结构化证据三项组件的移除结果。完整模型保留三项组件；其余设置仅移除名称所示组件，并以完整模型为表内比较参照。由于该表中的完整模型数值与主实验表不同，消融结果用于固定协议下的相对比较，不与表~\ref{tab:main-results} 的绝对数值直接比较。

\begin{landscape}
\begin{table}[H]
\centering
\caption{三个公开教育数据集上的关键组件消融实验结果}
\label{tab:ablation_study}
\small
\setlength{\tabcolsep}{5pt}
\resizebox{\linewidth}{!}{
\begin{tabular}{lcccccccccccc}
\toprule
\multirow{2}{*}{\textbf{组件设置}}
& \multicolumn{4}{c}{\textbf{FoundationalASSIST}}
& \multicolumn{4}{c}{\textbf{XES3G5M}}
& \multicolumn{4}{c}{\textbf{MOOCRadar}} \\
\cmidrule(lr){2-5}
\cmidrule(lr){6-9}
\cmidrule(lr){10-13}
& AUC & ACC & F1 & RMSE
& AUC & ACC & F1 & RMSE
& AUC & ACC & F1 & RMSE \\
\midrule
完整模型
& \textbf{0.7022} & \textbf{0.6830} & 0.7602 & \textbf{0.4972}
& \textbf{0.6627} & 0.7834 & 0.8745 & \textbf{0.4338}
& \textbf{0.7500} & 0.8172 & \textbf{0.8976} & 0.4062 \\

移除目标条件认知证据检索
& 0.6532 & 0.6409 & 0.7302 & 0.5170
& 0.6120 & 0.7320 & 0.8366 & 0.4606
& 0.5974 & 0.7234 & 0.8303 & 0.4660 \\

移除结构化认知摘要
& 0.6990 & 0.6810 & 0.7668 & 0.4997
& 0.6571 & \textbf{0.7837} & 0.8750 & 0.4363
& 0.7375 & \textbf{0.8178} & 0.8964 & \textbf{0.3983} \\

移除结构化证据
& 0.6784 & 0.6600 & \textbf{0.7703} & 0.5311
& 0.6375 & 0.7815 & \textbf{0.8752} & 0.4367
& 0.6873 & 0.8115 & 0.8938 & 0.4041 \\
\bottomrule
\end{tabular}
}
\end{table}
\end{landscape}

\begin{figure}[H]
\centering
\includesvg[inkscapelatex=false,width=0.92\textwidth]{figures/fig3_ablation_auc_effects_v2}
\caption{关键组件消融的 AUC 差值。每个点表示完整模型与对应移除组件变体之间的 \(\Delta\)AUC；误差线表示配对 bootstrap 95\% 置信区间。}
\label{fig:ablation-auc-effects}
\end{figure}

图~\ref{fig:ablation-auc-effects} 进一步给出了完整模型相对各移除变体的 AUC 差值及配对 bootstrap 置信区间。从 AUC 指标看，移除目标条件认知证据检索后，FoundationalASSIST、XES3G5M 和 MOOCRadar 上的完整模型 AUC 相对该变体分别高出 0.0458、0.0451 和 0.1526，且 95\% 置信区间均为正。该结果表明，在当前消融设置下，目标条件认知证据检索是三个数据集上最稳定的概率排序贡献来源，尤其在 MOOCRadar 上作用更明显。

移除结构化证据同样会带来稳定的 AUC 下降，三个数据集的差值分别为 0.0207、0.0196 和 0.0627，说明证据角色、支持强度和关系来源等结构化信息为目标条件预测提供补充。相比之下，移除结构化认知摘要后的 AUC 下降幅度较小，说明摘要信息对概率排序有增益，但弱于目标条件认知证据检索和结构化证据。

这表明完整模型的主要收益并非来自让 LLM 自行总结或生成判断，而是来自目标条件认知证据检索和结构化证据组织；结构化认知摘要更适合解释为对检索证据的辅助压缩。需要注意的是，部分阈值指标与误差指标在个别变体上存在方向不完全一致的情况。因此，本文将消融结论限定为“目标条件认知证据检索对 AUC 最稳定、结构化证据提供次级支持、结构化认知摘要提供较小补充”，不将任何组件解释为对所有指标均一致有效。

考虑到消融实验需要重复执行证据构造、检索与冻结 LLM 推理，本文在固定测试目标子集上开展关键组件比较。所有消融变体共享相同的训练数据与测试目标集合，以保证组件比较的可比性。主实验仍在完整测试集上报告，消融实验用于分析不同组件的相对贡献。


\section{结论}

本文围绕训练阶段未见知识概念上的归纳式知识追踪，提出了 Cognitive-RAG KT。该方法的核心不是扩大历史输入规模，而是以目标题目为条件，从学生目标前历史中选择可迁移认知证据，并将分层认知上下文、结构化认知摘要派生特征和目标题目侧先验映射为 soft slots 注入冻结 LLM，最终由标签 logits 输出作答正确概率。

实验结果表明，Cognitive-RAG KT 在未见概念归纳评估下优于多类知识追踪基线；纯未见概念切片说明其表现不是主要由混合已见概念样本驱动；关键组件消融进一步表明，目标条件认知证据检索是最稳定的贡献来源，结构化证据提供次级支持，结构化认知摘要提供辅助压缩信息。本文的主要结论是：对于训练阶段未见概念上的知识追踪，显式选择和组织学生目标前历史中的可迁移认知证据，比无条件压缩历史序列更能为冻结 LLM 的概率判别提供有效信息。

\section{局限性与未来工作}

尽管 Cognitive-RAG KT 在三个公开数据集上取得了稳定的整体提升，但本文仍存在若干需要进一步研究的边界。第一，当前实验通过整体比较、切片分析、分层分析和组件消融支持了目标条件认知证据检索的有效性，但尚未充分刻画被检索证据的认知可信性。性能增益并不必然等价于证据在教学认知层面具有充分性、必要性或因果有效性。未来可引入证据遮蔽、反事实替换、检索扰动、人工专家评价和证据一致性分析，检验模型依赖的历史交互是否真正支撑当前预测，而不是利用题目相似性、概念频率或数据集偏差获得表面相关性。

第二，结构化认知摘要是对检索证据的受约束压缩，而不是经过独立验证的真实认知诊断。本文仅将其作为学生侧上下文特征和摘要派生特征使用，并通过消融评估其预测贡献；摘要字段本身是否与教师判断、学生错因或后续学习收益一致，仍需要人工标注、专家评价或干预实验进一步验证。

第三，本文将知识追踪任务限定为目标前历史条件下的作答正确概率预测，并主要采用 AUC、ACC、F1、RMSE 和 BCE 等离线指标评价。该设置能够检验模型的预测能力，但尚不足以说明模型输出能否直接服务于教学干预。

后续研究应进一步考察预测概率的校准性、低频概念与新学生场景上的稳定性，以及模型输出与学习收益、错因诊断和个性化干预之间的关系。对于基于 LLM 的知识追踪方法，仅报告离线预测提升仍不充分；更关键的问题是模型依赖和呈现的认知证据是否可验证、可追溯且对教师有用。

第四，冻结 LLM 增强了结构化认知上下文的条件表示与概率判别能力，但也显著提高了推理成本，使其在大规模在线教育平台中的直接部署仍面临限制。未来可探索证据缓存、soft slots 蒸馏、轻量化学生模型替代、检索候选预计算以及检索器—预测器联合压缩等策略。

同时，后续工作需要系统分析检索规模、slot 数量、LLM 参数规模和推理延迟之间的性能—成本权衡。只有在保持预测性能和证据可信性的同时降低部署成本，Cognitive-RAG KT 才能进一步从离线实验框架走向真实教育系统应用。





\clearpage
\bibliographystyle{plainnat}
\bibliography{references}

\clearpage
\appendix
\section{附录}
\label{app:experiment-details}

本附录补充主文未展开的评估协议、主模型置信区间、分层结果范围、对比模型设置、方法配置、消融设计与计算环境。为保持主文紧凑，附录仅保留有助于复现和审阅的关键信息；脚本入口、字段与中间清单只在消融协议等必要位置列出。

\subsection{评估协议补充说明}

本文的未见概念归纳评估以题目关联概念为划分依据。所有测试目标均至少包含一个训练阶段未见概念；依赖学生作答记录构建的题目难度、协同关系和图结构统计仅由训练交互估计；测试标签仅用于评价指标计算，不参与题目先验构建、证据选择或模型选择。

\begin{table}[htbp]
\centering
\small
\caption{未见概念评估协议补充说明}
\label{tab:appendix-audit-summary}
\resizebox{\textwidth}{!}{%
\begin{tabular}{llll}
\toprule
\textbf{数据集} & \textbf{测试目标约束} & \textbf{行为统计来源} & \textbf{目标标签用途} \\
\midrule
FoundationalASSIST & 至少包含一个训练阶段未见概念 & 仅由训练交互估计 & 仅用于评价指标计算 \\
XES3G5M & 至少包含一个训练阶段未见概念 & 仅由训练交互估计 & 仅用于评价指标计算 \\
MOOCRadar & 至少包含一个训练阶段未见概念 & 仅由训练交互估计 & 仅用于评价指标计算 \\
\bottomrule
\end{tabular}%
}
\end{table}

\subsection{主模型置信区间与分层分析补充结果}

\begin{table}[htbp]
\centering
\small
\caption{Cognitive-RAG KT 主模型指标及 bootstrap 置信区间}
\label{tab:appendix-main-ci}
\resizebox{\textwidth}{!}{%
\begin{tabular}{lrrrrr}
\toprule
\textbf{数据集} & \textbf{样本数} & \textbf{AUC} & \textbf{ACC} & \textbf{F1} & \textbf{RMSE} \\
\midrule
FoundationalASSIST & 666,034 & 0.7805 [0.7793, 0.7816] & 0.7335 [0.7324, 0.7345] & 0.7998 [0.7989, 0.8007] & 0.4239 [0.4233, 0.4245] \\
XES3G5M & 496,715 & 0.7542 [0.7525, 0.7559] & 0.8046 [0.8035, 0.8058] & 0.8873 [0.8866, 0.8880] & 0.3791 [0.3781, 0.3800] \\
MOOCRadar & 818,197 & 0.8404 [0.8393, 0.8416] & 0.8521 [0.8513, 0.8528] & 0.9111 [0.9107, 0.9116] & 0.3322 [0.3314, 0.3328] \\
\bottomrule
\end{tabular}%
}
\end{table}

\begin{figure}[htbp]
\centering
\includesvg[inkscapelatex=false,width=0.98\textwidth]{figures/figs1_stratified_auc_ranges_v2}
\caption{不同分层维度下的 AUC 最小--最大范围。单元格颜色表示同一分层维度内各 bucket 的 AUC 跨度，单元格文字给出最小值和最大值；该范围为描述性分层范围，不表示置信区间。}
\label{fig:stratified-auc-ranges}
\end{figure}

图~\ref{fig:stratified-auc-ranges} 中的范围为同一分层维度内各 bucket 的最小值和最大值，用于描述模型在不同测试条件下的性能变化，不表示置信区间。

\subsection{对比模型的输入信息}

\begin{table}[htbp]
\centering
\small
\caption{对比模型的类别与输入信息}
\label{tab:appendix-comparison-models}
\begin{tabularx}{\textwidth}{p{2.6cm}p{3.3cm}X}
\toprule
\textbf{模型} & \textbf{类型} & \textbf{输入信息} \\
\midrule
DKT & 循环神经网络 KT & 学生历史题目、知识概念与作答序列；不使用额外认知上下文 \\
SAKT & 自注意力 KT & 学生历史交互序列与当前目标题目；不使用额外认知上下文 \\
SAINT & Transformer KT & 题目序列与作答序列；不使用额外认知上下文 \\
AKT & 注意力与认知测量 KT & 学生历史交互序列、题目或概念表示；不使用额外认知上下文 \\
DIMKT & 难度感知 KT & 学生历史交互序列和题目难度相关信息；不使用 Cognitive-RAG 证据 \\
QIKT & 问题级 KT & 学生历史交互序列和题目级表示；不使用 Cognitive-RAG 证据 \\
TCKT & 时序/因果增强 KT & 学生历史交互序列及其时序关系；不使用额外认知上下文 \\
SimpleKT & 简化注意力 KT & 学生历史交互序列；不使用额外认知上下文 \\
SparseKT & 稀疏注意力 KT & 学生历史交互序列；不使用额外认知上下文 \\
DenoiseKT & 去噪 KT & 学生历史交互序列及去噪机制；不使用 Cognitive-RAG 证据 \\
RobustKT & 鲁棒 KT & 学生历史交互序列；不使用额外认知上下文 \\
KeenKT & 近期知识追踪模型 & 学生历史交互序列；不使用额外认知上下文 \\
Cognitive-RAG KT & 本文方法 & 学生历史交互、目标条件认知证据、结构化认知摘要、题目侧认知先验和 soft slots \\
\bottomrule
\end{tabularx}
\end{table}

\subsection{Cognitive-RAG KT 模型配置}

\begin{table}[htbp]
\centering
\small
\caption{Cognitive-RAG KT 的关键模型配置}
\label{tab:appendix-main-method-config}
\begin{tabularx}{\textwidth}{p{3.8cm}X}
\toprule
\textbf{项目} & \textbf{设置} \\
\midrule
语言模型骨干 & Qwen3-8B，主体参数冻结 \\
预测方式 & 读取标签 logits，经二分类 softmax 得到 \(\hat{p}_{u,t}\) \\
soft-slot 模式与数量 & context-target；学生侧 context soft slots 数量为 4，目标题目侧 target soft slots 数量为 2 \\
学生侧输入特征 & 目标条件上下文嵌入、结构化认知摘要嵌入、摘要派生结构特征和学生侧数值状态特征 \\
目标题目侧输入特征 & 题目文本语义向量、分层语义 ID 向量、认知增强语义向量、协同题目向量 \\
结构化认知摘要 schema & 输出为单个 JSON object，顶层字段固定为 \texttt{mastered\_concepts}、\texttt{weak\_concepts}、\texttt{transfer\_state}、\texttt{risk\_level}、\texttt{evidence\_quality}、\texttt{diagnosis}。其中 \texttt{mastered\_concepts} 和 \texttt{weak\_concepts} 最多各 3 项，每项包含 \texttt{concept}、\texttt{evidence\_ids} 和 \texttt{confidence}；风险、证据质量和置信度均限定为低、中、高。 \\
摘要生成约束 & 摘要提示要求只输出 JSON，不输出解释、推理过程或 Markdown；实现使用兼容 chat-completion 的 JSON object 响应格式、温度 0.0 和最大输出长度 256 tokens，并对字段完整性、枚举值、证据编号和诊断长度进行解析校验。不合法输出最多触发三次重试；仍不合法时不写入摘要字段。 \\
摘要编码方式 & 合法摘要被转写为稳定掌握、薄弱风险、诊断说明和辅助状态等短文本通道，并形成摘要侧嵌入；同时构造 11 维 \texttt{llm\_struct\_features}，包括掌握点数量、薄弱点数量、掌握/薄弱存在指示、诊断长度归一化、风险等级 one-hot 和证据质量 one-hot。摘要嵌入与结构特征合并到学生侧 context projector 的输入中。 \\
投影器设置 & 隐藏维度为 512，dropout 为 0.1 \\
原始长文本输入 & 默认不将完整历史长文本直接输入 LLM，仅使用经过投影的 soft slots \\
检索与重排设置 & 历史窗口 20，第一阶段候选数 30，最终证据数 6，重排候选数 30，重排权重为 0.5 \\
训练设置 & 固定随机种子 42；冻结 LLM，仅训练 soft-slot 概率预测模块中的投影模块 \\
\bottomrule
\end{tabularx}
\end{table}

\subsection{对比模型训练超参数}

\begin{table}[htbp]
\centering
\small
\caption{对比模型训练协议与模型超参数}
\label{tab:appendix-baseline-training-protocol}
\resizebox{\textwidth}{!}{%
\begin{tabular}{lrrrrl}
\toprule
\textbf{模型} & \textbf{训练轮数} & \textbf{批大小} & \textbf{学习率} & \textbf{最大序列长度} & \textbf{主要结构参数} \\
\midrule
DKT & 100 & 256 & 0.001 & 100 & emb=100, hidden=100 \\
SAKT & 100 & 256 & 0.001 & 100 & d=100, heads=5, dropout=0.2 \\
SAINT & 100 & 256 & 0.001 & 100 & d=30, heads=5, dropout=0.6 \\
AKT & 100 & 256 & 0.001 & 100 & d\_model=100, heads=5, dropout=0.2 \\
DIMKT & 100 & 256 & 0.001 & 100 & emb=100, dropout=0.2 \\
QIKT & 100 & 256 & 0.001 & 100 & emb=100, dropout=0.2 \\
TCKT & 100 & 256 & 0.001 & 512 & d\_model=100, heads=5, dropout=0.2 \\
SimpleKT & 100 & 256 & 0.001 & 512 & d\_model=100, heads=5, dropout=0.2 \\
SparseKT & 100 & 256 & 0.001 & 512 & d\_model=100, heads=5, sparse\_ratio=0.8 \\
DenoiseKT & 100 & 256 & 0.001 & 512 & d\_model=100, dropout=0.2, dropout1=0.1 \\
RobustKT & 100 & 256 & 0.001 & 100 & d\_model=100, heads=5, dropout=0.2 \\
KeenKT & 100 & 256 & 0.0001 & 512 & d\_model=256, heads=8, dropout=0.2 \\
\bottomrule
\end{tabular}%
}
\end{table}

\subsection{关键组件消融协议}

表~\ref{tab:appendix-ablation-design} 给出主文表~\ref{tab:ablation_study} 中各变体的定义。所有变体共享相同的训练数据、相同随机种子和同一固定测试目标子集；配对 bootstrap 分析在每个数据集 50,000 个目标样本上进行。主实验仍在完整测试集上报告，因此消融表只用于固定协议下的组件相对比较。

固定测试目标子集由主实验测试预测文件导出为评估清单，生成入口为 \texttt{scripts/build\_llm\_ablation\_eval\_manifest.py}。清单中的每一行以三元组 \((u,t,q_{u,t})\) 对应的 \texttt{sample\_id} 作为唯一键，并记录数据集名、测试 split、学生 ID、目标时间步、目标题目 ID 和真实标签。

若指定抽样规模，则采用按标签分层的随机抽样，默认随机种子为 42；若抽样规模为 0，则保留全部测试行。清单生成后同时写出 metadata，记录源预测文件、源测试样本数、入选样本数、抽样策略、随机种子、正负标签数量和清单 SHA-256。

后续消融运行由 \texttt{scripts/run\_llm\_ablation\_experiments.py} 读取该固定清单，不再使用新的 \texttt{sample\_size}、\texttt{max\_samples} 或 offset 参数。运行配置记录 \texttt{eval\_manifest\_jsonl}、各变体名称、加载样本数和入选标签分布，并检查清单标签与学生数据中的目标标签一致。该协议保证表~\ref{tab:ablation_study} 中不同变体的差异来自被移除组件，而不是测试目标或标签组成变化。

\begin{table}[htbp]
\centering
\small
\caption{关键组件消融设置}
\label{tab:appendix-ablation-design}
\begin{tabularx}{\textwidth}{p{3.3cm}XXX}
\toprule
\textbf{变体} & \textbf{目标条件认知证据检索} & \textbf{结构化认知摘要} & \textbf{结构化证据} \\
\midrule
完整模型 & 保留 & 保留 & 保留 \\
移除目标条件认知证据检索 & 移除 & 保留 & 保留 \\
移除结构化认知摘要 & 保留 & 移除 & 保留 \\
移除结构化证据 & 保留 & 保留 & 移除 \\
\bottomrule
\end{tabularx}
\end{table}

\subsection{计算环境}

\begin{table}[htbp]
\centering
\small
\caption{实验计算环境}
\label{tab:appendix-experiment-environment}
\begin{tabular}{ll}
\toprule
\textbf{项目} & \textbf{配置} \\
\midrule
GPU & 6 $\times$ NVIDIA GeForce RTX 4090 \\
CPU 与内存 & Intel Xeon Platinum 8350C，503 GiB RAM \\
操作系统 & Ubuntu 22.04.5 LTS \\
主要软件环境 & Python 3.11.15，PyTorch 2.6.0+cu124，Transformers 4.51.1 \\
\bottomrule
\end{tabular}
\end{table}


\end{document}
